\section{Investment analysis, conflicts and responsibilities as a member}

\begin{qfrule}[Applying Standard V, investment analysis, recommendations and actions: diligence and reasonable basis, communication with clients, and record retention, with conforming and violating cases]\label{rule:standard-investment-analysis}
\qffield{Statement} \emph{V(A) Diligence and reasonable basis.} Analyse investments with diligence, independence and thoroughness, and have a reasonable and adequate basis, supported by research, for every recommendation or action. The depth depends on the product and the member's role. Secondary and third-party research may be relied on after reasonable checks of its assumptions, rigour, timeliness and objectivity, or in good faith where the firm has vetted the provider. Quantitative models must be understood and tested, including scenarios outside the historical data; their developers need more diligence than their users. External advisers are selected and reviewed as diligently as securities. A member need not dissociate from a group recommendation he or she disagrees with if it has a reasonable basis \citep[pp.~129--134]{standards-practice-handbook-2024}.

\emph{V(B) Communication with clients and prospective clients.} Disclose (1) the nature of services and their costs to the client, including costs of affiliates and third parties, kept up to date; (2) the basic format and principles of the investment process and material changes, including the use of external advisers; (3) significant risks and limitations, such as leverage, liquidity and capacity; (4) the factors important to the analysis; and (5) distinguish fact from opinion. Clients' sophistication does not excuse missing disclosures, though they may be tailored (pp.~141--146).

\emph{V(C) Record retention.} Keep records that support analyses, recommendations, actions and client communications, including e-mails, messages and social-media posts. Records belong to the firm; without regulatory or firm rules, keep them at least seven years (pp.~153--154).

\qffield{Practice} Contrasting cases (original).
\begin{itemize}
\item \emph{Backtest.} A quant team launches a factor strategy after a backtest on 2010--2021 data without testing a rate shock or transaction costs. \emph{Violation} of V(A): the model was not tested outside the observed regime. Launching after stress tests, out-of-sample validation and capacity analysis, documented, \emph{conforms} even if the strategy later loses.
\item \emph{Third-party research.} A small adviser relies on a large broker's reports after reviewing the broker's methods and its record. \emph{Conforming}. Relying on an anonymous blog's earnings forecast without checking \emph{violates} V(A).
\item \emph{Costs.} An adviser's agreement lists her 0.80 percent advisory fee but not that the recommended funds' share classes pay her firm 0.25 percent. \emph{Violation} of V(B)(1); the affiliated revenue is part of the client's cost.
\item \emph{Fact and opinion.} A report states ``earnings will rise 15 percent next year'' as if certain. \emph{Violation} of V(B)(5); ``we estimate'' with the assumptions shown conforms.
\item \emph{Records.} A departing analyst copies her models to use at a new firm. \emph{Violation} of V(C) and IV(A): she must rebuild the records at the new firm from sources she may lawfully use.
\end{itemize}

\qffield{Conditions} Handbook, 12th edition; the cost disclosure in V(B)(1) took effect in 2024.

\qffield{Pitfalls} Treating a loss as proof of insufficient diligence, or a gain as proof of enough. Disclosing the process once and never again.

\qffield{See also} \hyperref[meth:preventing-analysis-violations]{prevention}; \hyperref[meth:backtest-design-evaluation]{backtest design}.

\end{qfrule}

\begin{qfmethod}[Preventing violations of Standard V: research-review and model-testing policies, criteria for outside providers and external managers, written service and cost disclosures, and retention schedules]\label{meth:preventing-analysis-violations}
\qffield{Statement} \citep[pp.~133--134, 146, 154]{standards-practice-handbook-2024}:
\begin{itemize}
\item policies requiring that research, ratings and recommendations have a substantiated basis, with written due-diligence procedures;
\item measurable criteria for the quality of research and the accuracy of recommendations over time;
\item written minimum scenario testing for all models used to develop, rate or evaluate instruments, covering the breadth of scenarios, accuracy over time and sensitivity of cash flows to inputs;
\item criteria and review frequency for outside data and research providers and for external advisers and sub-advisers;
\item a rigorous pre-publication review of research, and records of the research behind each report;
\item written disclosure of services and costs at the outset, updated when they change (best practice, p.~142);
\item archived research notes and communications even where the firm's own retention policy is narrower.
\end{itemize}

\qffield{Practice} \emph{Model governance at Aster Quant} (original): no strategy goes live without (1) a research log recording every specification tried, (2) a holdout period never used in development, (3) cost and capacity estimates, (4) stress scenarios including 2008, 2020 and a 300 basis point rate rise, and (5) sign-off by a reviewer who did not build the model. Client agreements state the fee, fund expenses, any affiliated revenue and an annual estimate in currency of total costs. Every recommendation is filed with its inputs and outputs for seven years. Items 1--4 answer the backtest case above; the agreement answers the cost case.

\qffield{Conditions} Recommendations; regulators may prescribe model-risk and retention rules.

\qffield{Pitfalls} A model review by the model's author. Cost disclosure that omits third-party fund expenses.

\qffield{See also} \hyperref[rule:standard-investment-analysis]{Standard V}; \hyperref[meth:strategy-promotion-criteria]{strategy promotion}.

\end{qfmethod}

\begin{qfrule}[Applying Standard VI, conflicts of interest: avoid or disclose conflicts, priority of transactions, and referral fees, with conforming and violating cases]\label{rule:standard-conflicts-of-interest}
\qffield{Statement} \emph{VI(A) Avoid or disclose conflicts.} A conflict is any matter that could reasonably be expected to impair independence and objectivity or interfere with duties to clients, prospective clients or the employer. Avoid it where reasonable; otherwise disclose it prominently, in plain language and effectively, and mitigate it. Typical conflicts: ownership of recommended securities (beneficial ownership includes family holdings and voting or dispositive power), board seats, investment-banking, underwriting and market-making relationships, fee arrangements and fund rebates, and incentive pay \citep[pp.~157--160]{standards-practice-handbook-2024}.

\emph{VI(B) Priority of transactions.} Client and employer transactions come before those in which the member is a beneficial owner. Personal investing is allowed if clients are not disadvantaged, the member does not benefit from client trades, and rules are followed. Personal transactions include those for immediate family members' accounts and any account in which the member has a direct or indirect financial interest. Family accounts that are client accounts are treated like any other firm account, neither given special treatment nor disadvantaged because of the family relationship. That holds even when the member is a beneficial owner of the family client account; the beneficial ownership then brings the firm's preclearance and reporting requirements for personal transactions, not a lower priority (pp.~164--165; Example~3, pp.~167--168).

\emph{VI(C) Referral fees.} Disclose to the employer, clients and prospective clients, before engagement, any compensation or benefit received or paid for recommending products or services: its nature (flat or percentage, one-off or continuing, cash, soft dollars or in kind) and estimated value (p.~169).

\qffield{Practice} Contrasting cases (original).
\begin{itemize}
\item \emph{Inheritance.} An analyst's spouse inherits USD~2 million of a stock he covers. \emph{At minimum} he must disclose it to his employer and in his reports; best practice is to hand coverage to a colleague. Continuing without disclosure \emph{violates} VI(A).
\item \emph{Boilerplate.} A firm's brochure says it ``may receive compensation from product sponsors'' when it receives 0.25 percent a year from every fund on its platform. \emph{Inadequate}: the conflict exists and must be described specifically (compare the SEC's 2019 interpretation, which rejects ``may'' for existing conflicts, \citep[84 FR 33676]{sec-adviser-standard-of-conduct-2019}).
\item \emph{Front-running.} A portfolio manager buys a stock for herself an hour before placing a large client order in it. \emph{Violation} of VI(B). Buying it two days after the client orders are complete, with pre-clearance, \emph{conforms}.
\item \emph{Personal preference.} A manager has brokers put part of each scarce IPO into her husband's brokerage account, which the firm does not manage, while her clients' orders go unfilled. \emph{Violation} of VI(B): clients must be filled first, and the account must be disclosed to the employer.
\item \emph{Family client treated fairly.} A manager's parents are fee-paying discretionary clients in the firm's balanced strategy. In an oversubscribed IPO suitable for that strategy, their account gets the same pro-rata fill as the other accounts and is flagged for compliance review. This \emph{conforms}. Filling them only after every other client, to avoid any accusation of favouritism, \emph{violates} VI(B).
\item \emph{Family client with beneficial ownership.} A manager is joint owner, with her mother, of a fee-paying taxable investment account run in the same strategy on standard terms. The account's participation in the block is precleared, the account is on her holdings report and compliance receives its statements; it gets the same pro-rata fill as other suitable accounts. This \emph{conforms}. Giving it a smaller fill or none because of her interest disadvantages a client (VI(B), III(B)); skipping the preclearance and reporting the firm requires breaches the firm's code that VI(B) expects her to follow (p.~166). Topping it up from other clients' fills would put her interest first and \emph{violates} VI(B).
\item \emph{Selling against advice.} A manager sells a stock she still recommends because she needs cash for a house, after clients' trades. \emph{No violation}: personal needs can justify trades counter to advice.
\item \emph{Referral.} An adviser pays an accountant 10 percent of first-year fees for referrals and tells new clients only after they sign. \emph{Violation} of VI(C): disclosure must precede the agreement.
\end{itemize}

\qffield{Conditions} The ``avoid'' option was added in 2024 (Handbook, p.~ix).

\qffield{Pitfalls} Treating disclosure as a cure for every conflict. Excluding family holdings from beneficial ownership. Letting beneficial ownership decide a family client account's allocation: ownership triggers preclearance and reporting, management decides the allocation.

\qffield{See also} \hyperref[meth:preventing-conflict-violations]{prevention}; \hyperref[def:broker-agency-conflicts]{broker conflicts}.

\end{qfrule}

\begin{qfmethod}[Preventing violations of Standard VI: conflict registers and plain-language disclosure, pre-clearance, blackout periods, IPO and private-placement limits, holdings reporting and referral-fee approval]\label{meth:preventing-conflict-violations}
\qffield{Statement} \citep[pp.~157--160, 165--167, 169]{standards-practice-handbook-2024}:
\begin{itemize}
\item a register of conflicts, reviewed when their nature changes, with disclosures updated and written in plain language, prominent rather than in footnotes;
\item firewalls for staff who sit on issuer boards;
\item personal-investing policies: limited or no participation in IPOs, pre-approval of private placements, blackout periods around client trades, disclosure of holdings at hiring and periodically, duplicate confirmations and statements, pre-clearance of trades, and disclosure of the policy itself to clients in specific terms;
\item referral-fee procedures: prohibit them or require approval, and report the amounts received at least quarterly to the employer.
\end{itemize}

\qffield{Practice} \emph{Personal-trading code for Brook Lane Advisers} (original): (1) pre-clear all trades in securities held or under review for clients; (2) a seven-day blackout before and after client trades in the same security; (3) a 30-day minimum holding period; (4) no IPOs for personal accounts, meaning employees' own accounts and other accounts in which they or immediate family have a beneficial interest that are not managed as client accounts; relatives' fee-paying accounts managed as client accounts are allocated as clients even when the employee has a beneficial interest, but those accounts remain subject to items (1) and (5), with preclearance confirming eligibility and never the size of the pro-rata fill, while participation in the client block is exempt from the blackout in item (2); private placements only with written approval; (5) duplicate statements to compliance; (6) monthly comparison of employee and client trades; (7) clients receive the code. In the front-running case the pre-clearance would have been refused; in the referral case, the policy requires a written disclosure before engagement giving the 10 percent rate and an estimated amount.

\qffield{Conditions} Recommendations; US advisers are separately required to adopt codes of ethics covering personal trading.

\qffield{Pitfalls} Pre-clearance without post-trade comparison. A conflict register that is not shown to clients.

\qffield{See also} \hyperref[rule:standard-conflicts-of-interest]{Standard VI}.

\end{qfmethod}

\begin{qfrule}[Applying Standard VII, responsibilities as a CFA Institute member or candidate: conduct in programmes and references to the designation, with conforming and violating cases]\label{rule:standard-member-responsibilities}
\qffield{Statement} \emph{VII(A) Conduct as participants in CFA Institute programmes.} Do not compromise the reputation or integrity of CFA Institute or the designation or the integrity, validity or security of its programmes: no cheating, breaching testing rules, disclosing exam content (including topics or formulas tested or not tested), circumventing security, misusing an association with CFA Institute or misrepresenting information to it. Volunteers must not disclose confidential material. Members remain free to criticise CFA Institute's policies \citep[pp.~173--175]{standards-practice-handbook-2024}.

\emph{VII(B) Reference to CFA Institute, the designation and the programme.} Do not misrepresent or exaggerate what membership, the designation or candidacy means. Charterholders must be current members, having filed the annual Professional Conduct Statement and paid dues; lapsed members may say only that they were members. Candidates must say they are candidates and must not imply a partial designation or an expected completion date. Factual statements are allowed; claims of superior performance are not (pp.~180--183).

\qffield{Practice} Contrasting cases (original).
\begin{itemize}
\item \emph{Exam content.} After the exam, a candidate posts that ``there was a lot on swaps and nothing on GIPS''. \emph{Violation} of VII(A). Posting that the exam was hard and she found the time pressure difficult is \emph{allowed}.
\item \emph{Designation.} A résumé says ``CFA Level III''. \emph{Violation} of VII(B): there is no partial designation. ``Passed Level III of the CFA Program; eligible for the charter after the work requirement'' \emph{conforms}.
\item \emph{Lapsed membership.} A charterholder stops paying dues but keeps ``CFA'' on her LinkedIn profile. \emph{Violation}; she must remove it until reinstated.
\item \emph{Marketing.} ``Our CFA charterholders achieve better results.'' \emph{Violation}. ``Our analysts hold the CFA charter, which requires passing three examinations and committing annually to the Code and Standards'' \emph{conforms}.
\end{itemize}

\qffield{Conditions} Applies to members and candidates; policies for each programme set further rules.

\qffield{Pitfalls} Treating informal comments about an exam's emphasis as harmless. Treating the designation as a lifetime title irrespective of membership.

\qffield{See also} \hyperref[meth:preventing-member-violations]{prevention}; \hyperref[def:professional-conduct-programme]{enforcement}.

\end{qfrule}

\begin{qfmethod}[Preventing violations of Standard VII: guidance and templates for references to the designation, membership checks and exam-confidentiality training]\label{meth:preventing-member-violations}
\qffield{Statement} Recommended practices: circulate the guidance on Standard~VII(B) to the firm's legal, compliance, public-relations and marketing teams, and use templates for biographies and marketing that reference the designation correctly \citep[p.~183]{standards-practice-handbook-2024}. Candidates should study the testing policies before an exam and refrain from discussing its content in any forum (pp.~173--174). Firms can check annually that staff described as charterholders are current members (Interpretation, from the membership requirements, pp.~180--181).

\qffield{Practice} \emph{Annual designation check at Kerr \& Lowe} (original): compliance reconciles the list of staff named as charterholders on the website, pitch books and e-mail signatures with CFA Institute's public member directory each January; biographies use three approved forms (charterholder; candidate in the CFA Program; passed Level~II); new joiners who are candidates receive a one-page note on exam confidentiality before their first exam season. The lapsed-membership and ``CFA Level III'' cases above would be caught by the first and second steps.

\qffield{Conditions} The directory check is an interpretation of how to apply the membership requirements, not a Handbook procedure.

\qffield{Pitfalls} Marketing teams editing biographies without compliance review.

\qffield{See also} \hyperref[rule:standard-member-responsibilities]{Standard VII}.

\end{qfmethod}
